\gccFrontHeading{摘要}

本文围绕高校校园网建设场景，设计并规划了一套面向教学、科研与日常管理业务的校园网络方案。系统以数据通信技术为基础，结合核心层、汇聚层与接入层的分层架构，完成网络拓扑设计、地址规划、路由策略、安全控制与运维管理等关键内容。方案重点关注现代高校网络中常见的高并发访问、跨区域互联、资源访问控制和后期扩展问题，力求在稳定性、安全性与可维护性之间取得平衡。

在研究过程中，本文首先分析校园网建设的业务需求与约束条件，明确教学区、办公区、宿舍区和公共服务区的网络接入特点；随后基于主流网络设备与协议，提出核心交换、冗余链路、虚拟局域网划分、访问控制列表和网络地址转换等设计方法；最后从安全、兼容性、可扩展性和运行效率等方面对方案进行评价。设计结果表明，该方案能够较好地支撑高校校园网络的基本业务运行，并为后续数字校园建设提供网络基础。

\vspace{0.6em}
\noindent\textbf{关键词：}\gccKeywordsZh

